\documentclass[11pt]{article}

\usepackage[margin=0.9in]{geometry}
\usepackage{amsmath,amssymb}
\usepackage{booktabs}
\usepackage{enumitem}
\usepackage{listings}
\usepackage{xcolor}
\usepackage{hyperref}
\usepackage{tikz}

\definecolor{backcolour}{rgb}{0.95,0.95,0.92}
\definecolor{codeblue}{rgb}{0.05,0.15,0.55}
\definecolor{codegreen}{rgb}{0.0,0.45,0.0}

\lstset{
    backgroundcolor=\color{backcolour},
    basicstyle=\ttfamily\small,
    keywordstyle=\color{codeblue}\bfseries,
    commentstyle=\color{codegreen},
    stringstyle=\color{purple},
    breaklines=true,
    frame=single,
    showstringspaces=false,
    tabsize=4
}

\title{CPSC 250 \\ Inheritance and Polymorphism}
\author{Edward Brash}
\date{}

\begin{document}

\maketitle

\section{Motivation}

In many programs, several classes share many of the same features.

If we write those features separately in every class, we duplicate code.

That is inefficient and hard to maintain.

Inheritance gives us a better strategy.

\section{The Basic Idea}

Inheritance allows us to define a base class containing common information and behavior.

Derived classes then inherit from that base class and add specialized information.

\section{Example: Businesses}

Suppose we want to model several kinds of businesses.

A general \texttt{Business} class might contain:

\begin{itemize}
    \item name,
    \item address,
    \item hours of operation.
\end{itemize}

A \texttt{Restaurant} class might include all of that plus:

\begin{itemize}
    \item cuisine type,
    \item rating.
\end{itemize}

A \texttt{CarDealership} class might include all business information plus:

\begin{itemize}
    \item car makes sold,
    \item average price by make.
\end{itemize}

\section{Example: Grocery Store Items}

Another example is grocery store items.

The base class might be:

\[
\texttt{Item}
\]

with:

\begin{itemize}
    \item name,
    \item quantity.
\end{itemize}

Derived classes might include:

\begin{itemize}
    \item \texttt{Produce},
    \item \texttt{Cereal},
    \item \texttt{Meat}.
\end{itemize}

Each derived class has the base features, plus additional specialized features.

\section{Base Class: Item}

\begin{lstlisting}[language=Python]
class Item:

    def __init__(self):
        self.name = ""
        self.quantity = 0

    def set_name(self, name):
        self.name = name

    def set_quantity(self, quantity):
        self.quantity = quantity

    def display(self):
        print(self.name, self.quantity)
\end{lstlisting}

\section{Derived Class: Produce}

\begin{lstlisting}[language=Python]
class Produce(Item):

    def __init__(self):
        Item.__init__(self)
        self.expiration = ""

    def set_expiration(self, expiration):
        self.expiration = expiration

    def get_expiration(self):
        return self.expiration
\end{lstlisting}

\section{Calling the Base Class Constructor}

The line:

\begin{lstlisting}[language=Python]
Item.__init__(self)
\end{lstlisting}

is an explicit call to the base class constructor.

This initializes the inherited variables:

\begin{itemize}
    \item \texttt{name},
    \item \texttt{quantity}.
\end{itemize}

This must be done because Python does not automatically call the base constructor in this style.

\section{Main Program Example}

\begin{lstlisting}[language=Python]
item1 = Item()
item1.set_name("Cereal")
item1.set_quantity(9)
item1.display()

item2 = Produce()
item2.set_name("Apples")
item2.set_quantity(40)
item2.set_expiration("5/5/2024")
item2.display()

print(f"Expires on {item2.get_expiration()}")
\end{lstlisting}

\section{Why This Is Useful}

The derived class \texttt{Produce} can use inherited methods from \texttt{Item}:

\begin{itemize}
    \item \texttt{set\_name()},
    \item \texttt{set\_quantity()},
    \item \texttt{display()}.
\end{itemize}

It only needs to define what is new:

\begin{itemize}
    \item expiration date.
\end{itemize}

\section{Polymorphism}

Polymorphism means that different classes can be treated in a common way.

For example, if several objects all have a method called:

\begin{lstlisting}[language=Python]
display()
\end{lstlisting}

then code can call \texttt{display()} without needing to know exactly which type of object it is handling.

\section{Inheritance Chains}

Inheritance can continue over multiple levels.

For example:

\[
\texttt{Base}
\rightarrow
\texttt{DerivedOne}
\rightarrow
\texttt{DerivedTwo}
\]

Each class adds information to the previous class.

\begin{lstlisting}[language=Python]
class Base:

    def __init__(self):
        self.name = ""


class DerivedOne(Base):

    def __init__(self):
        Base.__init__(self)
        self.address = ""


class DerivedTwo(DerivedOne):

    def __init__(self):
        DerivedOne.__init__(self)
        self.zip_code = ""
\end{lstlisting}

\section{Important Design Point}

The notes emphasize that explicitly calling the base class constructor may feel unusual if you come from other languages.

However, it reflects the spirit of inheritance:

\begin{quote}
The derived class should intentionally initialize the base-class portion of itself.
\end{quote}

\section{Key Takeaways}

\begin{itemize}
    \item Inheritance reduces duplicated code.
    \item A base class stores common features.
    \item A derived class adds specialized features.
    \item Derived classes can use inherited methods.
    \item Python often requires explicit base constructor calls.
    \item Polymorphism allows different objects to be used through common methods.
\end{itemize}

\end{document}
